\textbf{Duración:} 3 horas (en aula)\\
\textbf{Enfoque:} Distinguir chat de sistemas agénticos; human-in-the-loop; integridad académica en contextos EIC (ecuaciones, citas, datos).

\subsubsection*{Objetivos de la sesión}
Al finalizar, la persona participante:
\begin{enumerate}
\item Explica con ejemplos propios la diferencia entre un chat ad hoc y un flujo agéntico.
\item Identifica roles típicos (planificador, redactor, revisor, verificador) y puntos HITL.
\item Formula una nota de integridad (qué automatiza, qué decide la persona) ligada a material STEM.
\item Propone un tema EIC tentativo para el microartefacto del CADi.
\end{enumerate}

\subsubsection*{Agenda (resumen)}
\begin{tabularx}{\textwidth}{@{}lclX@{}}
\toprule
Bloque & Tiempo & Contenido \\
\midrule
Apertura y mapa 40 h & 15 min & Duración 24+16; evidencias; expectativas \\
Conceptos & 30 min & Agentes, herramientas, memoria; preview orquestación \\
Integridad y riesgos & 25 min & Alucinación, sesgo, datos; ecuaciones y citas \\
Demo & 25 min & Flujo con HITL; anti-patrones \\
Taller & 50 min & Mapa + nota de integridad + tema EIC \\
Cierre & 15 min & Puesta en común; TI lecturas (bloque A) \\
\bottomrule
\end{tabularx}

\subsubsection*{Producto esperado}
Mapa conceptual agentes vs.\ chat + nota de integridad (5--8 líneas) + tema EIC tentativo (E1).

\subsubsection*{Trabajo independiente relacionado}
Sesión 2: brief, prompt de sistema y orquestación. TI: $\sim$1\,h lecturas/demos (bloque A).

\textit{Deck Beamer canónico:} \texttt{tex/beamer/sesion-01.tex}.
